\documentclass{article}
\usepackage{amsmath, amssymb, amsthm}
\usepackage{hyperref}
\usepackage{geometry}
\geometry{margin=1in}

\title{Erd\H{o}s Problem \#398}
\date{Last updated: 2025-08-31}
\author{Source: erdosproblems.com}

\begin{document}
\maketitle

\noindent\textbf{Status:} FALSIFIABLE \quad \textbf{No prize}

\noindent\textbf{Tags:} number theory, factorials

\medskip
\noindent\textbf{Note:} Brocard-Ramanujan conjecture

\medskip

\noindent\textbf{Problem Statement:}

Are the only solutions to\[n!=x^2-1\]when $n=4,5,7$?




The Brocard-Ramanujan conjecture. Erd\H{o}s and Graham describe this as an old conjecture, and write it 'is almost certainly true but it is intractable at present'. Overholt \cite{Ov93} has shown that this has only finitely many solutions assuming a weak form of the ABC conjecture . There are no other solutions below $10^9$ (see the OEIS page ). Naciri \cite{Na25} has proved that there are only finitely many solutions if $x\pm 1$ is either $k$-free (for some $k\geq 2$) or a prime power (and Cambie explains both facts in the comments), and that if $x\pm 1$ is $7$-free then $n=4,5,7$ give the only solutions.




References

[Na25] A. M. Naciri, On the Brocard-Ramanujan equation with $7$-free integers and prime powers . Integers (2025), A71.

[Ov93] Overholt, Marius, The Diophantine equation $n!+1=m^2$ . Bull. London Math. Soc. (1993), 104.

\noindent\textbf{Related OEIS sequences:} A146968, A141399

\bigskip
\noindent\small{Source: \url{https://www.erdosproblems.com/398}}
\end{document}
